% Aula 07 --- sob o LDA o log-odds é linear; a recíproca é falsa.
Renata é atuária numa seguradora e quer classificar os segurados de um ramo em dois
perfis, $Y=1$ (risco alto) e $Y=0$ (risco baixo), a partir de uma única covariável: o
tempo $X$, em anos, entre a contratação da apólice e o primeiro sinistro. Ela hesita
entre a análise discriminante linear (LDA) e a regressão logística, dois métodos que
chegam a uma fronteira linear por caminhos diferentes. Ao longo do exercício,
$\sigma(u)=e^u/(1+e^u)$ é a função logística, $\log$ é o logaritmo natural e
$\pi_c=\Prob(Y=c)$.
\begin{itens}
  \item \pts{2,2} Comece pelo caso geral, com $p$ covariáveis: suponha que $\X\in\R^p$ e
        que, como no LDA, $\X\mid Y=c\sim N(\bm\mu_c,\bm\Sigma)$ para $c\in\{0,1\}$, com
        $\bm\Sigma$ simétrica e positiva definida. Mostre que
        $\Prob(Y=1\mid\X=\x)=\sigma(\beta_0+\bbeta^\top\x)$ e escreva $\beta_0$ e
        $\bbeta$ em função de $\pi_0$, $\pi_1$, $\bm\mu_0$, $\bm\mu_1$ e $\bm\Sigma$.
        Conclua que, sob as hipóteses do LDA, o modelo da regressão logística vale
        exatamente. (Escreva $\Prob(Y=1\mid\X=\x)$ pelo Teorema de Bayes, com $f_c$ a
        densidade de $\X$ dado $Y=c$, divida o numerador e o denominador por
        $\pi_1f_1(\x)$ e ponha o resultado na forma $\sigma\big(u(\x)\big)$. Depois,
        substitua as densidades normais em $u(\x)$, expanda as duas formas quadráticas e
        observe o que se cancela porque a covariância é a mesma nas duas classes.)
  \item \pts{1,8} Os tempos de Renata não se parecem com uma normal: são positivos e
        muito assimétricos. Suponha agora que $X\mid Y=c$ tenha distribuição
        exponencial de taxa $\lambda_c$, com densidade $\lambda_c e^{-\lambda_c x}$ para
        $x>0$, com $\lambda_0=1$, $\lambda_1=3$ e $\pi_0=\pi_1=1/2$: em média, o
        primeiro sinistro do perfil de risco alto chega três vezes mais cedo. Mostre que
        $\Prob(Y=1\mid X=x)=\sigma(\log3-2x)$ para $x>0$. O modelo da regressão
        logística vale nesta população, isto é, existem constantes $\beta_0$ e $\beta_1$
        com $\Prob(Y=1\mid X=x)=\sigma(\beta_0+\beta_1x)$ para todo $x>0$? As hipóteses do
        LDA, de que $X$ é normal em cada classe e tem a mesma variância nas duas, valem?
\end{itens}

\begin{solucao}
\textbf{(a)} Pelo Teorema de Bayes, com $f_c$ a densidade de $\X$ dado $Y=c$,
\[
  \Prob(Y=1\mid\X=\x)
  =\frac{\pi_1f_1(\x)}{\pi_0f_0(\x)+\pi_1f_1(\x)}
  =\frac{1}{1+e^{-u(\x)}}=\sigma\big(u(\x)\big),
\]
com
\[
  u(\x)=\log\frac{\pi_1f_1(\x)}{\pi_0f_0(\x)} ,
\]
em que a segunda igualdade sai de dividir numerador e denominador por $\pi_1f_1(\x)>0$.
Basta, portanto, mostrar que o \emph{log-odds} $u(\x)$ é afim em $\x$. As duas densidades
normais têm a mesma constante de normalização, $(2\pi)^{-p/2}(\det\bm\Sigma)^{-1/2}$,
porque a covariância é comum, e ela se cancela no quociente:
\[
  u(\x)=\log\frac{\pi_1}{\pi_0}
        -\tfrac12(\x-\bm\mu_1)^\top\bm\Sigma^{-1}(\x-\bm\mu_1)
        +\tfrac12(\x-\bm\mu_0)^\top\bm\Sigma^{-1}(\x-\bm\mu_0).
\]
Como $\bm\Sigma^{-1}$ é simétrica,
$(\x-\bm\mu)^\top\bm\Sigma^{-1}(\x-\bm\mu)
=\x^\top\bm\Sigma^{-1}\x-2\bm\mu^\top\bm\Sigma^{-1}\x+\bm\mu^\top\bm\Sigma^{-1}\bm\mu$. O
termo quadrático $\x^\top\bm\Sigma^{-1}\x$ aparece nas duas formas e se cancela, e sobra
\[
  u(\x)=\beta_0+\bbeta^\top\x,\qquad
  \bbeta=\bm\Sigma^{-1}(\bm\mu_1-\bm\mu_0),\qquad
  \beta_0=\log\frac{\pi_1}{\pi_0}
     -\tfrac12\big(\bm\mu_1^\top\bm\Sigma^{-1}\bm\mu_1-\bm\mu_0^\top\bm\Sigma^{-1}\bm\mu_0\big).
\]
Logo $\Prob(Y=1\mid\X=\x)=\sigma(\beta_0+\bbeta^\top\x)$, que é exatamente o modelo da
regressão logística. A covariância comum foi usada duas vezes, para cancelar as constantes de
normalização e para cancelar o termo quadrático, e só o segundo uso é indispensável para que
o \emph{log-odds} seja afim. Com covariâncias $\bm\Sigma_0\ne\bm\Sigma_1$, as constantes
deixariam em $u(\x)$ a parcela $\tfrac12\log(\det\bm\Sigma_0/\det\bm\Sigma_1)$, que não
depende de $\x$ e iria para $\beta_0$; já o termo
$-\tfrac12\,\x^\top\big(\bm\Sigma_1^{-1}-\bm\Sigma_0^{-1}\big)\x$ não seria identicamente nulo,
e o \emph{log-odds} seria quadrático em $\x$, como no QDA.

\smallskip
\textbf{(b)} Pela mesma fórmula de Bayes, para $x>0$,
\[
  u(x)=\log\frac{\pi_1\lambda_1e^{-\lambda_1x}}{\pi_0\lambda_0e^{-\lambda_0x}}
      =\log\frac{\pi_1}{\pi_0}+\log\frac{\lambda_1}{\lambda_0}-(\lambda_1-\lambda_0)\,x
      =0+\log3-2x ,
\]
e portanto $\Prob(Y=1\mid X=x)=\sigma(\log3-2x)$. O modelo logístico vale exatamente,
com $\beta_0=\log3\approx1{,}0986$ e $\beta_1=-2$. As hipóteses do LDA falham duas vezes:
$X\mid Y=c$ não é normal (vive em $(0,\infty)$ e é assimétrica), e as variâncias das duas
classes, $1$ e $1/9$, são diferentes. O item (a) mostra que as hipóteses do LDA bastam
para o modelo logístico, e este item mostra que não são necessárias: a logística supõe
apenas que o \emph{log-odds} seja afim, e isso acontece aqui sem nada de normal.
\end{solucao}
